\chapter{Additional Notes}

\section{Universes and Type Families}\label{universes-and-type-families}

The definition of a universe has an apparent circularity: a universe is defined using the concept of a type family, while type families are formalized as functions into a universe. The circularity is resolved by distinguishing between the \textsl{meta-theoretic} (judgmental) framework and the \textsl{object-theoretic} (internal) language of type theory.

\subsection*{Judgmental Type Families}

Before the introduction of any universes, the type theory relies on a primitive notion of type dependence. At this stage, a type family is not a mathematical object or a function that can be evaluated or composed: it is a syntactic rule, a \textsl{judgment}. It states that under the assumption of a variable $x$ of type $A$, an expression $P(x)$ forms a valid type. This is denoted using the notation
\[
    x\colon A\vdash P(x)\text{ type}.
\]
These judgmental families are rigidly bound to the deductive rules of the system and cannot be treated as first-class entities.

\subsection*{Internal Type Families}

In order to be able to manipulate type families as mathematical objects (e.g., quantify over them, compose them, and map into them), the judgmental families need to be internalized as functions. A universe $\univ$ is introduced to serve as the codomain for these functions
\[
    P\colon A\to\univ.
\]
Under this approach, the elements of $\univ$ are typically viewed as \textsl{codes} (or \textsl{names}) for types, rather than types themselves.

\subsection*{Tarski-Style Universes}

To define the universe $\univ$ without circularity, we rely on the pre-existing judgmental framework.

\begin{defn}
    A \textsl{Tarski-style universe} consists of a type $\univ$ of \textsl{codes} together with a decoding function
    \[
        X\colon\univ\vdash\tau(X)\text{ type}
    \]
    that assigns a type to each code. To accurately reflect the ambient type theory, the universe must be closed under type constructors. This means there exist code-forming operators in $\univ$ which $\tau$ decodes definitionally into the standard type formations:
    \begin{enumerate}[a),font=\upshape]
        \item \textsc{Dependent functions:} For any code $a\colon\univ$ and any family of codes $b\colon\tau(a)\to\univ$, there is a code $\hat{\Pi}(a,b)\colon\univ$ such that
        \[
            \tau(\hat{\Pi}(a,b))\equiv\prod_{x\colon\tau(a)}\tau(b(x)).
        \]
        \item \textsc{Dependent pairs:} For any code $a\colon\univ$ and family of codes $b\colon\tau(a)\to\univ$, there is a code $\hat{\Sigma}(a,b)\colon\univ$ such that
        \[
            \tau(\hat{\Sigma}(a,b))\equiv\sum_{x\colon\tau(a)}\tau(b(x)).
        \]
        \item \textsc{Identity types:} For any code $a\colon\univ$ and terms\/ $x,y\colon\tau(a)$, there is a code $\hat\Id(a,x,y)\colon\univ$ such that
        \[
            \tau(\hat\Id(a,x,y))\equiv x\eq{\tau(a)}y.
        \]
        \item \textsc{Coproducts:} For any codes $a,b\colon\univ$, there is a code $\mathop{\hat{+}}(a,b)\colon\univ$ such that
        \[
            \tau(\mathop{\hat{+}}(a,b))\equiv\tau(a)+\tau(b).
        \]
    \end{enumerate}
\end{defn}

\begin{rem}
    The assertion that the universe $\univ$ is itself a type cannot be derived from the decoding function $\tau$. Instead, its existence must be postulated via a primitive formation rule:
    \[
        \vdash \univ \text{ type}.
    \]
    One might naturally wonder whether $\univ$ is an element of itself, i.e., whether the judgment $\univ\colon\univ$ holds. However, positing that a universe contains its own code is notoriously inconsistent. Much like the set of all sets in classical set theory, it leads to a type-theoretic contradiction known as \textsl{Girard's paradox}.
    
    To circumvent this inconsistency while maintaining the ability to compute with universes, Homotopy Type Theory introduces an infinite hierarchy of universes:
    \[
        \univ_0, \univ_1, \univ_2, \dots
    \]
    where each universe is a term of the next: $\univ_i\colon\univ_{i+1}$. In a strict Tarski-style presentation, each universe $\univ_i$ has its own decoding function $\tau_i\colon\univ_i\to\univ_{i+1}$.
    
    In mathematical practice, managing these indices explicitly is exceedingly cumbersome. Following a convention known as \textsl{typical ambiguity}, we routinely drop the subscripts and simply write $\univ$, implicitly trusting that a consistent assignment of universe levels could be reconstructed to avoid paradoxes.
\end{rem}

\subsection*{Russell-Style Universes}

While the Tarski-style separation between codes and types is necessary for foundational rigor, managing the decoding function $\tau$ explicitly becomes notationally burdensome. In mathematical practice, it is customary to adopt a \textsl{Russell-style} presentation. This is achieved by systematically omitting the decoding function, identifying the code $A\colon\univ$ with its corresponding type $\tau(A)$.

Under this convention, the universe $\univ$ is treated as if it were a type whose elements are types. The decoding properties of $\tau$ ensure that this abuse of notation is perfectly safe, as the code-forming operators (such as $\hat{\Pi}$ and $\hat{\Sigma}$) become completely transparent and indistinguishable from the standard type constructors. 

\begin{ntn}
    With the Russell-style convention in place, an internal type family is simply a function
    \[
        P\colon B\to\univ.
    \]
    For any element $b\colon B$, the application $P(b)$ is an element of $\univ$, which we directly regard as a type. Consequently, we can immediately form the dependent types
    \[
        \prod_{b\colon B}P(b) \qquad\text{and}\qquad \sum_{b\colon B}P(b)
    \]
    without writing $\tau(P(b)$.
\end{ntn}

\subsection*{Types by Computation}

A profound consequence of having a universe is the ability to define types by \textsl{computation}. Because $\univ$ is a valid codomain for functions, we can construct type families using the elimination principles (pattern matching and induction) of other types. This is fundamentally impossible with purely judgmental families.

\begin{xmpl}
    We can use a computed family over the boolean type $\fun2$ to formally prove that its constructors $1_{\fun2}$ and $0_{\fun2}$ are distinct. We define a function $P\colon\fun2\to\univ$ by pattern matching:
    \begin{align*}
        P(1_{\fun2})&\defeq\fun1\\
        P(0_{\fun2})&\defeq\fun0.
    \end{align*}
    If there were a path $p\colon1_{\fun2}\eq{\fun2}0_{\fun2}$, we could transport the canonical element $\star\colon\fun1$ along $p$ in the family $P$. This transport would yield an element of the empty type:
    \[
        \fun{transport}^P(p,\star)\colon\fun0.
    \]
    Since $\fun0$ has no elements, such a path $p$ cannot exist. This standard derivation, which shows $1_{\fun2}\neq0_{\fun2}$, relies essentially on the existence of the universe $\univ$: since $\univ$ is a type, we can use it as the codomain in the induction principle for~$\fun2$.
\end{xmpl}

\section{Yoneda's Lema}

\begin{thm}\label{thm:yoneda-general}
    {\upshape[Yoneda’s Lemma]} Let\/ $A$ be an object in a category\/~$\cat C$. Let\/ $F\colon\cat C\to\cat{Set}$ be a covariant functor, and let\/ $\yoneda_A=\Hom_{\cat C}(A,-)$ be the representable functor from\/ $\cat C$ to\/ $\cat{Set}$. Then the set\/ $F(A)$ is in a one-to-one correspondence with the collection of natural transformations from\/~$\yoneda_A$ to\/~$F$. This correspondence maps each natural transformation\/ $\eta\colon\yoneda_A\to F$ to the element\/ $\eta_A(\id_A)$ in\/~$F(A)$.
\end{thm}

\begin{proof}
    Let $c\in F(A)$ be an element. For any object $Y$ define
    \begin{align*}
        \eta_{c,Y}\colon \yoneda_A(Y)&\to F(Y)\\
        \zeta&\mapsto F(\zeta)(c).
    \end{align*}
    Given $u\colon Y\to Z$ we have $\yoneda_A(u)(\zeta)=u\circ\zeta$ for $\zeta\in\yoneda_A(Y)$. Put $(\eta_c)_Y=\eta_{c,Y}$. Then $\eta_c\colon\yoneda_A\to F$ is natural because the diagram
    $$
        \begin{tikzcd}
            \yoneda_A(Y)
                    \arrow[r,"(\eta_c)_Y"]
                    \arrow[d,"\yoneda_A(u)"']
                &F(Y)
                    \arrow[d,"F(u)"]
                &\zeta
                    \arrow[r,mapsto]
                    \arrow[d,mapsto]
                &F(\zeta)(c)
                    \arrow[d,mapsto]\\
            \yoneda_A(Z)
                    \arrow[r,"(\eta_c)_Z"']
                &F(Z)&u\circ\zeta
                    \arrow[r,mapsto]
                &F(u\circ\zeta)(c)
        \end{tikzcd}
    $$
    commutes (since $F(u)(F(\zeta)(c))=F(u\circ\zeta)(c)$ by the functoriality of $F$).
    
    Conversely, take a natural transformation $\eta\colon\yoneda_A\to F$. Since $\eta_A$ takes endomorphisms of $A$ to elements of $F(A)$, we know that $c_\eta=\eta_A(\id_A)$ is such an element.

    It remains to be seen that both maps are reciprocal.

    \begin{description}
        \item[$\mbf{F(A)\to\op{Nat}\to F(A):}$] Take $c\in F(A)$. Then,
        \begin{align*}
            c&\mapsto\eta_c\mapsto c_{\eta_c}\\
            c_{\eta_c} &= (\eta_c)_A(\id_A)\\
                &= \eta_{c,A}(\id_A)\\
                &= F(\id_A)(c)\\
                &= \id_{F(A)}(c)\\
                &= c.
        \end{align*}
        \item[$\mbf{\op{Nat}\to F(A)\to\op{Nat}:}$] Take $\eta\colon\yoneda_A\to F$. Given $\zeta\colon A\to Y$, we have a commutative diagram
        $$
            \begin{tikzcd}
                \yoneda_A(A)
                        \arrow[d,"\yoneda_A(\zeta)"']
                        \arrow[r,"\eta_A"]
                    &F(A)
                        \arrow[d,"F(\zeta)"]
                    &\id_A
                        \arrow[r,mapsto]
                        \arrow[d,mapsto]
                    &c_\eta
                        \arrow[d,mapsto]\\
                \yoneda_A(Y)
                        \arrow[r,"\eta_Y"']
                    &F(Y)&\zeta\circ\id_A=\zeta
                        \arrow[r,mapsto]
                    &F(\zeta)(c_\eta)
            \end{tikzcd}
        $$
        Therefore,
        \begin{align*}
            (\eta_{c_\eta})_Y(\zeta) &= \eta_{c_\eta,Y}(\zeta)\\
                &= F(\zeta)(c_\eta)\\
                &= \eta_Y(\zeta),
        \end{align*}
        which shows that $\eta_{c_\eta}=\eta$, as desired.
    \end{description}
\end{proof}

\section{The Slice Category}\label{slice-category}

\begin{defn}
    Given a set $B$ in $\cat{Set}$, the \textsl{slice category over $B$} is the category $\cat{Set}_B$ whose objects are the function $p\colon E\to B$, and its morphisms $\Hom_{\cat{Set}_B}(p\colon E\to B,p'\colon E'\to B)$ are the functions $\phi\colon E\to E'$ that make commutative the triangle
    \begin{equation}\label{tik:slice-morphism}
        \begin{tikzcd}[column sep=small]
            E
                    \arrow[rd,"p"']
                    \arrow[rr,"\phi"]
                &&E'
                    \arrow[ld,"p'"]\\
                &B
        \end{tikzcd}
    \end{equation}
\end{defn}

\begin{rem}
    Recall that an object $\mbf1$ in a category $\cat{C}$ is \textsl{terminal} if, for any other object $X$, there is a unique morphism $X\to\mbf1$. If such an object exists, a \textsl{global point} in an object $X$ is an arrow $s\colon\mbf1\to X$.
\end{rem}

\begin{prop}
    The identity\/ $\id_B\colon B\to B$ is a terminal object in\/ $\cat{Set}_B$. In particular, the global points of\/ $p$ are the sections of\/~$p$.
\end{prop}

\begin{proof}
    Let $p\colon E\to B$ be an arbitrary object of $\cat{Set}_B$. Then $p\colon E\to B$ is a morphism from $p$ to $\id_B$:
    \[
        \begin{tikzcd}[column sep=small]
            E
                    \arrow[rd,"p"']
                    \arrow[rr,"p",dashed]
                &&B
                    \arrow[ll,bend right=45,"s"',dashed]
                    \arrow[ld,"\id_B"]\\
                &B
        \end{tikzcd}
    \]
    and any morphism $s$ from $\id_B$ to $p$ is, by definition, of section of $p$.
\end{proof}

\begin{rem}
    Consider the set $B$ as a so called \textsl{discrete category} whose objects are its elements and its morphisms reduce to formal identities $\id_b\colon b\to b$ for $b\in B$. Given two functors $F,G\colon B\to\cat{Set}$, since the only morphisms in $B$ are the identities, a natural transformation $\eta\colon F\to G$ consists of a family of maps $\eta_b\colon F(b)\to G(b)$.
\end{rem}

\begin{ntn}
    As usual, we will denote the functor category from $B$ to $\cat{Set}$ by $\fun{Fun}(B,\cat{Set})$.
\end{ntn}

\begin{thm}\label{thm:grothendieck-construction}
    {\upshape[Grothendieck Construction]} Let\/ $B$ be a set. Then, there is a natural isomorphism\/ $\cat{Set}_B\simeq\fun{Fun}(B,\cat{Set})$.    
\end{thm}

\begin{proof}
    The forward functor
    \[
        \mathcal F\colon\cat{Set}_B\to\fun{Fun}(B,\cat{Set})
    \]
    is defined on objects by
    \[\tag{$1$}
        \mathcal F(p)(b)=p^{-1}(b),
    \]
    for an object $p$ in $\cat{Set}_B$ and $b\in B$. For a morphism $\phi\colon p\to p'$ as in \eqref{tik:slice-morphism}, we define
    \[\tag{$2$}
        \mathcal F(\phi)_b\colon p^{-1}(b)\to (p')^{-1}(b),
        \quad x\mapsto\phi(x),
    \]
    for $b\in B$ and $x\in p^{-1}(b)$.

    The backward functor
    \[
        \mathcal G\colon\fun{Fun}(B,\cat{Set})\to\cat{Set}_B
    \]
    is defined on objects by the projection map
    \[\tag{$3$}
        \mathcal G(F)\colon\coprod_{b\in B}F(b)\to B,
    \]
    which maps each element of $F(b)$ to $b$ for every $b\in B$. For a natural transformation $\eta\colon F\to G$ between two functors in $\fun{Fun}(B,\cat{Set})$, we define the morphism
    \[\tag{$4$}
        \mathcal G(\eta)\colon\mathcal G(F)\to\mathcal G(G),
        \quad x\mapsto\eta_b(x),
    \]
    for $b\in B$ and $x\in F(b)$.

    To see that these functors are mutually inverse, observe that for any functor $F$ and $b\in B$,
    \begin{align*}
        (\mathcal F\circ\mathcal G)(F)(b)
            &=\mathcal F(\mathcal G(F))(b)\\
            &=\mathcal G(F)^{-1}(b)
                &&\text{; by }(1)\\
            &=F(b)
                &&\text{; by }(3).
    \intertext{For a natural transformation $\eta\colon F\to G$, evaluating at $x\in F(b)$ yields}
        (\mathcal F\circ\mathcal G)(\eta)_b(x)
            &=\mathcal G(\eta)(x)
                &&\text{; by }(2)\\
            &=\eta_b(x)
                &&\text{; by }(4),
    \end{align*}
    so $\mathcal F\circ\mathcal G$ is the identity functor on $\fun{Fun}(B,\cat{Set})$.

    Conversely, fix an object $p\colon E\to B$ in $\cat{Set}_B$. Then,
    \begin{align*}
        (\mathcal G\circ\mathcal F)(p)
            &=\mathcal G(\mathcal F(p))\\
            &=\coprod_{b\in B}p^{-1}(b)\xto\pi B\
                &&\text{; by }(1)\text{ \& }(3)\\
            &=E\xto pB\\
            &=p.
    \intertext{Now consider a morphism $\phi\colon p\to p'$. For $b\in B$ and $x\in p^{-1}(b)$, we have}
        (\mathcal G\circ\mathcal F)(\phi)(x)
            &=\mathcal G(\mathcal F(\phi))(x)\\
            &=\mathcal F(\phi)_b(x)
                &&\text{; by }(4)\\
            &=\phi(x)
                &&\text{; by }(2).
    \end{align*}
    Thus, $\mathcal G\circ\mathcal F=\Id$, the identity functor on $\cat{Set}_B$. This completes the proof.
    %
    
\end{proof}
    
Even if one can take some\/ $f\colon 1\to X$ and obtain the commutative diagram:
\[
    \begin{tikzcd}
        1 \arrow[d, "f"'] \arrow[dr, "y\circ f"] & \\
        X \arrow[r, "y"'] & Y,
    \end{tikzcd}
\]
there still may be other nonidentical arrows\/ $1\to X$, which provide nonequivalent `globalisations' of the local element\/ $y\in Y^X$. Hence,\/ $X$ can have non-global elements which are seen only from some stages of definition (`contexts of observation'). So, while the global elements of an object correspond to the set-theoretic idea of points, there can also be elements of an object that are not points. The importance of this idea cannot be overestimated. It provides a fresh view on many structures.